\chapter{Introducción}
\label{sec.c5:introduccion}

En esta actividad evaluativa correspondiente a la Unidad Temática 6 (Fuentes de Alimentación) se aborda el
cálculo, diseño y verificación experimental por simulación de una fuente de alimentación lineal
regulada con protección contra sobrecargas, en cumplimiento de los requerimientos fijados para la Consigna 5 del
Grupo 5.

\section{Especificaciones de diseño}
\label{sec.c5:especificaciones}

La consigna establece las siguientes condiciones operativas para el circuito:
\begin{itemize}
    \item Tensión de salida regulada: $V_{out} = \SI{12}{\volt}$ continua.
    \item Corriente máxima a plena carga: $I_{Lmax} = \SI{3}{\ampere}$.
    \item Corriente mínima de carga: $I_{Lmin} = \SI{100}{\milli\ampere}$.
    \item Componente regulador base: regulador lineal integrado de tres terminales LM7812 en encapsulado metálico TO-3.
    \item Protección: incorporación obligatoria de una etapa limitadora de sobrecarga y cortocircuito.
    \item Evaluación de desempeño: determinación de la totalidad de los factores de calidad de la fuente.
\end{itemize}

\section{Justificación de la topología seleccionada}
\label{sec.c5:topologia}

El circuito integrado LM7812 presenta una capacidad máxima de corriente nominal de salida de $\SI{1}{\ampere}$
(con un límite de disparo típico interno de $\SI{1.5}{\ampere}$ a $\SI{2.2}{\ampere}$). Dado que la carga
especificada demanda hasta $\SI{3}{\ampere}$, el integrado resulta insuficiente por sí solo para abastecer la
demanda continua sin ingresar en protección térmica o sobrecorriente interna.

Para superar esta restricción, adoptamos la configuración de \emph{elevador de corriente con transistor de paso
exterior y limitador de sobrecarga}, correspondiente a la Figura 6-67 del material teórico de la cátedra
\cite{gonzalez_fuentes3}. En este circuito, un transistor de potencia PNP ($Q_1$) conduce la mayor parte de la
corriente que excede la capacidad del regulador, mientras que un segundo transistor discreto ($Q_2$) junto a un resistor
de sensado ($R_{sc}$) forman un lazo de realimentación negativa para limitar rígidamente la corriente máxima
en caso de falla o cortocircuito franco en la salida.
